\documentclass[11pt]{article}
\usepackage[margin=1in]{geometry}
\usepackage{amsmath,amssymb}
\usepackage{iftex}
\ifPDFTeX
  \usepackage[T1]{fontenc}
  \usepackage[utf8]{inputenc}
\else
  \usepackage{unicode-math}
  \defaultfontfeatures{Scale=MatchLowercase}
\fi
\usepackage{lmodern}
\usepackage{microtype}
\usepackage{parskip}
\usepackage{array}
\usepackage[unicode,hidelinks]{hyperref}
\usepackage{bookmark}
\setlength{\emergencystretch}{3em}
\newcommand{\N}{\mathbb N}
\newcommand{\Th}{\operatorname{Th}}
\newcommand{\FV}{\operatorname{FV}}
\newcommand{\ucl}{\operatorname{ucl}}
\newcommand{\sk}{\operatorname{sk}}
\newcommand{\rk}{\operatorname{r}}
\newcommand{\supp}{\operatorname{supp}}
\newcommand{\cl}{\operatorname{cl}}
\newcommand{\MC}{\operatorname{MC}}
\newcommand{\Code}{\operatorname{Code}}
\newcommand{\cut}[2]{\Lambda_{#1}^{< #2}}
\newcommand{\eps}{\varepsilon}
\newcommand{\up}{\upsilon}
\newcommand{\qed}{\hfill\ensuremath{\square}}
\newenvironment{proof}{\par\emph{Proof.}\ }{\qed\par}
\newcommand{\statement}[2]{\par\textbf{#1.}\ #2\par}
\hypersetup{pdftitle={Mutual truthfulness and known mechanicalness: a proof of Alexander's Question 3-1-3},pdfauthor={}}
\title{Mutual truthfulness and known mechanicalness:\\
  a proof of Alexander's Question 3-1-3}
\author{}
\date{3 October 2026}

\begin{document}
\maketitle

\section*{Abstract}

Alexander's Question 3-1-3 asks whether a family of knowing machines can
know everyone's truthfulness, know particular codes for its subordinates,
and know everyone's mechanicalness. We give a proposed affirmative proof
for every recursively enumerable well-founded relation of subordination,
with no restriction on the formulas. A fixed common theory supplies
mechanicalness; the original machines have particular codes for this
theory and their direct subordinates. To prove the resulting family
sound, we use ordinal labels, auxiliary modal symbols that remember
their original owners, and finite hereditary restrictions of external
knowledge. Guarded transport recovers full source proof predicates at
prescribed labels. A separate core argument and an induction over whole
downward ideals then establish soundness for arbitrary hierarchies. We
extract the required effective ordinal structure from Wilken's pure
order-two resemblance analysis, including the domains of its partial
base conversions. The final reduct has precisely Alexander's original
operators. This paper presents the complete proposed hand argument;
its separate Lean formalization is unfinished. No literature priority
is claimed.

\section{The question and the affirmative statement}\label{sec:statement}

A machine may know that another machine has some program without
knowing which program it has. The distinction is the position of the
existential quantifier. For a formula $\phi$ with lone free variable $x$,
the two assertions are
\[
  \exists e\,K_i\forall x(K_j\phi\leftrightarrow x\in W_e),
  \qquad
  K_i\exists e\,\forall x(K_j\phi\leftrightarrow x\in W_e).
\]
The first supplies a particular index that $i$ knows to describe $j$'s
column. The second says that $i$ knows such an index exists. Here $W_e$
is the $e$th recursively enumerable subset of $\N$, and $x\in W_e$
abbreviates its standard arithmetical enumeration formula.

Let $\mathcal L$ be the language of Peano arithmetic with operators
$K_0,K_1,\ldots$. We use Alexander's weak modal semantics
\cite[\S\S1.2--1.3]{Alexander}. Knowledge is a formula-forming operator;
it is not given Kripke accessibility semantics. The arithmetic part of
our final structure is the standard natural numbers.

A family of theories $i\mapsto T_i$ determines knowledge by
\begin{equation}\label{eq:actual-semantics}
  M\models K_i\theta[s]\quad\Longleftrightarrow\quad
  T_i\vdash\theta^s.
\end{equation}
The sentence $\theta^s$ replaces each free variable by the numeral for
its value under $s$. Throughout, $\vdash$ means consequence in
Alexander's base logic. Completeness makes it equivalent to an
effective finite-certificate proof relation. A family of knowing
machines has recursively enumerable theories and satisfies the
self-doubting knowledge axioms: validity, deduction, introspection and
truthfulness, together with knowledge of the first three. Self-doubting
Epistemic Arithmetic additionally includes modal PA and knowledge of
its axioms. The induction schema ranges over the whole modal language.

Write $jRi$ when $j$ is subordinate to $i$. We assume $R$ is recursively
enumerable and well-founded. It need not be transitive. Let $R^+$ be its
positive transitive closure; a path of length zero is allowed when we
speak of a downward ideal.

\statement{Theorem 1 (full 3-1-3)}{For every recursively enumerable
well-founded relation $R$ on $\N$, there is a uniformly recursively
enumerable family of knowing machines satisfying self-doubting
Epistemic Arithmetic and the following schemas in $\mathcal L$:}
\begin{enumerate}
\item For all $i,j$ and every formula $\theta$,
  $K_i\ucl(K_j\theta\to\theta)$.
\item For $jRi$ and every $\phi$ with $\FV(\phi)=\{x\}$,
  $\exists e\,K_i\forall x(K_j\phi\leftrightarrow x\in W_e)$.
\item For all $i,j$ and every $\phi$ with $\FV(\phi)=\{x\}$,
  $K_i\exists e\,\forall x(K_j\phi\leftrightarrow x\in W_e)$.
\end{enumerate}

These are the three entries defining Question 3-1-3 in
\cite[Definition 1.1.1]{Alexander}. All formulas, including formulas
about superior and incomparable agents, are included in one family.
There is no bound on modal depth, arithmetical complexity or hierarchy
height. In particular, no recursive rank function for $R$ is assumed.

\statement{Corollary 2}{The same family gives affirmative answers to
Questions 3-1-1 and 3-1-2.}

Indeed, their mechanicalness requirements are subcollections of the
third schema. We therefore prove the stronger statement directly.

\subsection*{How the proof is organized}

The main construction is short. We add a temporary common knower $*$,
put all existence-of-code axioms in its theory $U$, and give each
original theory $T_i$ particular codes for $U$ and its subordinates.
Every $T_i$ can then inherit the closed theorems of $U$. What requires
work is proving that $U$ remains sound when all these code-holding
machines are present.

We first separate proof membership from truth by filtering each raw
proof predicate through the truth of its argument. This makes a
structure before soundness is known. Ordinal stages let us verify its
generators using smaller cuts. External modal occurrences that would
otherwise force a circular soundness assumption are restricted to
finitely many hereditary argument shapes. Auxiliary symbols record the
owners of discarded occurrences, so that arithmetic program indices
never have to change. The resulting local reconstruction permits an
induction on the original hierarchy. The two measures serve different
purposes: ordinal labels control active proof cuts; original hierarchy
ranks control the lower theories retained in a recursive call.

Sections~\ref{sec:fixed}--\ref{sec:masks} give the construction and its
local lemmas. Sections~\ref{sec:core}--\ref{sec:finish} establish the
global result. Section~\ref{sec:ordinals} proves the precise ordinal
input from published results and effective procedures, so the main
argument has no additional ordinal hypothesis.

\section{Fixed theories, fixed programs, and the truth filter}\label{sec:fixed}

Work temporarily in $\mathcal L^*=\mathcal L\cup\{K_*\}$. For a holder
$a\in\{*\}\cup\N$, let $B_a$ contain the following closed generators:
modal PA; Assigned Validity; $a$-Validity, $a$-Deduction and
$a$-Introspection; and $\ucl(K_b\theta\to\theta)$ for every $b,	heta$.
Assigned Validity means $\eta^s$ for every base-logically valid $\eta$
and numerical assignment $s$. All schematic formulas are universally
closed where appropriate.

If $\Gamma$ is a set of generating sentences, $\cl_a(\Gamma)$ contains
$K_a^n\sigma$ for $\sigma\in\Gamma$ and $n\ge0$, where the superscript
here denotes repeated application. We distinguish this ordinary finite
prefix from the ordinal-labelled heads introduced below. Write
$\Th(S)=\{\sigma:S\vdash\sigma\}$.

For a fresh index variable $e$, put
\[
  \MC_b(\phi)=\ucl\exists e\,\forall x
     (K_b\phi\leftrightarrow x\in W_e).
\]
Other free parameters of $\phi$ are closed \emph{outside} the
existential: an index may depend on their values. Define
\begin{equation}\label{eq:core}
 U=\cl_*\bigl(B_*\cup\{\MC_b(\phi):b,\phi\}\bigr).
\end{equation}
This recursively enumerable theory is fixed independently of $R$.
It has no particular-code generators. Choose an effective compiler
$d$ for all its unary consequence columns:
\begin{equation}\label{eq:d}
 n\in W_{d(\phi)}\quad\Longleftrightarrow\quad U\vdash\phi(\bar n),
 \qquad \FV(\phi)=\{x\}.
\end{equation}

Given a provisional family index $q$, let $c(q,j,\phi)$ be the program
enumerating the $\phi$-column of the theory with owner $j$ enumerated by
$q$. The parameter theorem makes $c$ effective. Define
\begin{equation}\label{eq:extensions}
\begin{split}
 T_i(q)=\cl_i\bigl(B_i
 &\cup\{\forall x(K_*\phi\leftrightarrow x\in W_{d(\phi)}):\FV(\phi)=\{x\}\}\\
 &\cup\{\forall x(K_j\phi\leftrightarrow x\in W_{c(q,j,\phi)}):
                   jRi,\ \FV(\phi)=\{x\}\}\bigr).
\end{split}
\end{equation}
The family is uniformly recursively enumerable in $q$: every edge is
admitted on a positive enumeration witness. The recursion theorem
chooses $q$ for this family. Set $T_i=T_i(q)$. Then
\begin{equation}\label{eq:c}
 n\in W_{c(q,j,\phi)}\quad\Longleftrightarrow\quad T_j\vdash\phi(\bar n).
\end{equation}
Equations~\eqref{eq:d} and \eqref{eq:c} are external identities of
enumerations. We never add an axiom saying that a theory knows its own
enumerator identity. The theories, $q$ and all displayed numerals are
fixed from now on.

\statement{Lemma 3 (own-known finite consequences)}{If
$S=\cl_a(\Gamma)$ contains own Validity and Deduction, then
$S\vdash\sigma$ for a sentence $\sigma$ implies $S\vdash K_a\sigma$.}

\begin{proof}
A consequence certificate uses finitely many generating premises
$\gamma_1,\ldots,\gamma_m$ and a valid implication from them to
$\sigma$. Their own prefixes belong to $S$. Own Validity gives knowledge
of the implication, and repeated own Deduction gives $K_a\sigma$.
The same argument iterates. It is a derived rule for sentences from
finite certificates, rather than an added global necessitation rule.
\end{proof}

For raw predicates $P_a$ on closed formulas, define $F_P$ recursively by
standard arithmetic and
\begin{equation}\label{eq:filtered}
 F_P\models K_a\theta[s]\quad\Longleftrightarrow\quad
 P_a(\theta^s)\ \text{and}\ F_P\models\theta[s].
\end{equation}
We call a raw family admissible when it is uniformly recursively
enumerable and invariant under alphabetic variation. We call $P_a$
\emph{fully sound in $F_P$} when $P_a(\sigma)$ implies
$F_P\models\sigma$ for every sentence $\sigma$. This refers to every
accepted sentence, not merely to assertions currently being inspected.

\statement{Lemma 4 (filtered structures)}{An admissible raw family
defines a legitimate Alexander structure $F_P$. Numerical assignment
commutes with truth: $F_P\models\theta[s]$ iff
$F_P\models\theta^s$. The structure satisfies modal PA and every
reflection instance $K_a\theta\to\theta$.}

\begin{proof}
The truth recursion in \eqref{eq:filtered} calls only a proper
subformula. Alphabetic invariance gives alphabetic invariance of truth.
Dependence only on free variables and the permitted variable-for-variable
substitution condition follow because numerical closure supplies the
same sentence on both sides, up to alphabetic variation. Formula
induction gives the assignment identity. Reflection is immediate from
the truth conjunct. For modal induction, the sentences at $x=\bar n$
and $x=S(\bar n)$ give the ordinary induction on $n$ in the standard
natural numbers. Here $S(\bar n)$ is literally $\overline{n+1}$.
This is Alexander's argument in Lemma 1.2.22 and Proposition 1.2.23,
and uses no arbitrary-term substitution law inside knowledge.
\end{proof}

The actual raw family will be $P_*=\Th(U)$ and $P_i=\Th(T_i)$.
Proving its full soundness removes every truth conjunct in
\eqref{eq:filtered} and yields \eqref{eq:actual-semantics}. Until then,
an enumerated proof alone does not verify a particular-code axiom.

\section{The ordinal room used by finite proofs}\label{sec:room}

We state the exact label facts now and derive them in
Section~\ref{sec:ordinals}. Let $L$ be a recursive notation domain for
the well-order $[0,\up_\omega)$, carrying recursive relations
$<,\le_1,\le_2$ from Wilken's \emph{pure} order-two structure.
Write $<_i$ for a strict comparison. The required properties are:
\begin{enumerate}
\item $\le_2\subseteq\le_1\subseteq\le$; the resemblance orders are
transitive. If $p\le_1r$ and $p<q\le r$, then $p\le_1q$.
\item If $\alpha\le_1\beta$, a finite expanded diagram below $\beta$
can be reflected below $\alpha$, fixing its already prescribed finite
part below $\alpha$. The language includes both resemblance relations,
and the diagram includes their negative facts.
\item A guard $G_k(\lambda)$ is witnessed by
$\lambda=g_0<_2g_1<_2\cdots<_2g_k$. If $k\ge1$, then above any finite
support $X<\lambda$ and any bound $u<\lambda$, there are arbitrarily
many finite copies $Y_0<\cdots<Y_{m-1}<\lambda$ of the first $k$ guard
points. Every copied point is $\le_1$-related to $\lambda$.
Thus $\max Y_t<\min Y_{t+1}$: separation concerns whole copies.
\item There is a cofinal strict chain
$C=\{\up_n:n\ge2\}$ with $\up_2<_2\up_3<_2\cdots$.
\end{enumerate}
The enumerators use finite relation checks and finite witnessed guards.
They need no decision procedure for membership in $C$. A $G_0$ guard
is one point. The room statement is only used with $k\ge1$.

There are two kinds of induction in this paper. The well-order $L$
supports effective labels for finite auxiliary proofs. The rank of $R$
will be used only as a metatheoretic measure; it never selects axioms.

\section{Recorded formulas and auxiliary presentations}\label{sec:recorded}

Fix a finite active set $A$ of holders and a finite set $E$ of old
external original names. There are two applications:
\begin{itemize}
\item In a \emph{core presentation}, $A=\{*\}$, the source holder is
$U$, and $E\subseteq\N$ is arbitrary. Main generators are the
mechanicalness generators of $U$.
\item In an \emph{extension presentation}, $A=J$ is a finite
$R^+$-antichain, the source holder at $a\in J$ is $T_a$, and $E$ is
disjoint from the whole original downward ideal $D_R(J)$. Main
generators are the particular core and subordinate codes of $T_a$.
There are no mechanicalness generators in this presentation.
\end{itemize}

For $e\in E$ introduce a fresh free modal symbol $Z_e$. An active head
may be labelled $K_a^\nu$, for $\nu\in L$. The \emph{source erasure}
$\eps$ removes ordinal labels, sends $Z_e$ to $K_e$, fixes every other
head and fixes arithmetic. Thus a ghost occurrence still remembers
its original owner. The adjective ``ghost'' describes the auxiliary
model; $Z_e$ is an ordinary unconstrained modal symbol in logical
certificates.

For example, an auxiliary version of a source code on $K_e\tau$ may
display $Z_e\tau$, but its program numeral is still the numeral for
the source argument $K_e\tau$. It is never the program for a different
original operator $K_s$. Erasure keeps the code identity while allowing
the auxiliary model to make the unwanted occurrence false.

Define the skeleton $\sk(\theta)$ by preserving formula constructors,
arithmetic predicate symbols and modal owners, while replacing terms
and binder variable identities by wildcards. Let $\rk(\theta)$ be
its constructor count. Proper formula subexpressions have smaller
rank. Alphabetic variation, numerical assignment and permitted variable
substitution preserve skeletons, including when two variables merge.

A finite hereditary protection table $H$ consists of pairs
$(e,\sk(\chi))$. The map $\pi_H$ keeps $K_e\chi$ live if its pair is
in $H$, and otherwise changes its head to $Z_e$; it applies recursively
to arguments. Heredity requires that a retained external argument
retain all its nested external arguments. Consequently
\begin{equation}\label{eq:erase-protect}
 \eps\pi_H\theta=\theta,
 \qquad\text{every retained external argument has a $Z$-free subtree.}
\end{equation}

A source-shared labelling assigns the same label to equal pairs
$(a,\sk(\eps\chi))$. Every active label lies strictly above
\emph{all} active labels anywhere in its argument, also through live
unindexed or ghost contexts. A formula satisfying this nesting is
globally admissible. Main mechanicalness and code images use one
protection table and one source-shared label table. Structural and
validity premises require global admissibility, without a requirement
that unrelated premises share a label table.

\statement{Lemma 5 (logical maps)}{Erasure, hereditary protection and
a global active stratifier preserve base-logical validity. They preserve
free variables, alphabetic variation and permitted substitutions, and
commute with numerical closure. Their finite certificate translations
are effective.}

\begin{proof}
In Alexander's first-order encoding $\Delta$, a modal atom $O\chi$ is
a separate predicate applied to the ordered tuple of its free
variables. The axioms BL1 and BL2 enforce alphabetic variation and
variable substitution. Each map sends this predicate to that of its
mapped modal atom, with the same free variables. Its images of BL1
and BL2 are again instances of those schemes. Renaming or identifying
predicates preserves a first-order validity certificate. Alternatively,
pull back an arbitrary target structure along the map and apply formula
induction; the substitution identities make the pullback a legitimate
weak-modal structure.

For a global stratifier, prescribe a consistent finite table first.
Process argument shapes in increasing rank, choosing each new label
above the finitely many labels of its descendants. Search in the
recursive unbounded order finds such a label. Repeated source shapes
receive the same entry. Thus the stratifier commutes with the required
substitutions. Numerical closure identities follow by recursion.

Only the stated direction of erasure is used. The converse would be
false: $K_e\tau\leftrightarrow Z_e\tau$ erases to a tautology but is
not a base-logical validity.
\end{proof}

\subsection*{The generator rules}

For each active holder $a$ enumerate a theory $\Lambda_a$. All
generators are closed sentences, with numerical instances admitted as
in the source. The following table specifies the structural part.
\begin{center}
\renewcommand{\arraystretch}{1.3}
\begin{tabular}{@{}p{.23\linewidth}p{.72\linewidth}@{}}
\textbf{Class}&\textbf{Admitted auxiliary generator}\\[3pt]
Arithmetic&Globally admissible modal PA, including full modal induction.\\
Assigned validity&$\eta^s$ for valid globally admissible $\eta$.\\
Own validity&$K_a^\nu\eta$ for valid $\eta$ with argument support below $\nu$.\\
Own deduction&Same-label deduction with every argument support below the label.\\
Own introspection&$K_a^\mu\eta\to K_a^\nu K_a^\mu\eta$, where $\mu<\nu$ and the whole formula is globally admissible.\\
Reflection&Truthfulness of indexed, unindexed and ghost heads.\\
Contextual collapse&Equivalence when one active label is changed along $\le_1$ inside a globally admissible context.\\
Own prefixes&Prefix any generating sentence at a holder label above its entire support, and iterate.\\
\end{tabular}
\end{center}
For contextual collapse, the changed argument lies below both labels;
all intermediate contexts used are globally admissible. It is available
to each holder. It is not an import of another holder's structural
axioms.

The main generators are labelled $\pi_H$-images of the source
mechanicalness sentences (core type) or particular codes (extension
type). In an extension image the displayed code is
\begin{equation}\label{eq:recorded-code}
 \forall x\bigl(K_w l\pi_H\phi
           \leftrightarrow x\in W_{p(w,\phi)}\bigr),
 \quad
 p(w,\phi)=
 \begin{cases}d(\phi),&w=*,\\c(q,w,\phi),&w\in\N.\end{cases}
\end{equation}
Its owner $w$ is neither active nor old external. Its numerator uses
the \emph{source} $\phi$, not the visible formula $l\pi_H\phi$.

Inspect the whole distinguished modal atom of a main image, including
its outer head. Whenever $K_e\chi$ is live with $e\in E$, every active
atom $K_b^\nu\beta$ inside $\chi$ must have a finite witness for
$G_{\rk(\eps\chi)}(\nu)$. Its argument support must be below $\nu$
and $\le_1$-related to $\nu$. Inspect nested live heads even below
ghosts. In particular, a core mechanicalness target $K_e\chi$ imposes
guards on the active atoms of its argument even when there is no other
external head in $\chi$.

Every finite guard point $h$ is made actual formula support by padding
the main sentence with an implication from holder-local tautology
markers $K_a^h\tau_h$. Choose distinct arithmetic tautology shapes
absent from the main formula. The tags are freely chosen formulas;
they do not encode the notation number of $h$. A reflected label keeps
the same tag. Each marker is provable by the holder's own validity.
The padded image therefore yields the unpadded main sentence by
discharging its markers.

All admission checks are finite and recursive on label notations.
Source validity and generator certificates are recursively enumerable.
For extensions, positive $R$-edge witnesses supply the subordinate
codes. The antichain and ideal-separation hypotheses apply to the
semantic uses of these theories; the enumeration makes no negative
$R$-test and consults no hierarchy rank.

\section{Erasure, collapse, and transport at a fixed label}\label{sec:transport}

Let $S_a$ be the source theory, either $U$ or $T_a$. The support of a
recorded formula is the finite set of its active labels, including
the marker labels. Let $\cut a\nu$ contain exactly the generated
sentences whose \emph{whole} support is below $\nu$, including those
with empty support. A finite consequence certificate contains its
premises, generator witnesses and a valid implication to its conclusion.
This follows from Alexander's compactness and completeness theorem
\cite[Theorem 1.2.6 and \S1.3]{Alexander}.

\statement{Lemma 6 (source erasure)}{If $\Lambda_a\vdash\eta$, then
$S_a\vdash\eps\eta$. The same holds for every cut.}

\begin{proof}
Each structural generator erases to its source counterpart; contextual
collapse erases to the equivalence of a formula with itself. A padded
main image erases to an implication from own-known arithmetic
tautologies to the identical source mechanicalness or code generator.
It is therefore a source consequence. A generated prefix erases to a
source consequence by Lemma 3, including prefixes of discharged code
images. Translate the final validity certificate by Lemma 5.
The source argument and its literal program numeral are unchanged.
\end{proof}

\statement{Lemma 7 (finite-support collapse)}{If
$\alpha\le_1\beta$ and $\supp(\eta)<\alpha$, then
\begin{equation}\label{eq:collapse}
 \cut a\alpha\vdash\eta\quad\Longleftrightarrow\quad
 \cut a\beta\vdash\eta.
\end{equation}}

\begin{proof}
The forward implication is inclusion. For the reverse, take the support
of a finite proof below $\beta$. Reflect its expanded finite diagram
below $\alpha$, fixing every support point already below $\alpha$,
in particular the target's labels. The reflection preserves strict
nesting, all guard witnesses, and each relation comparison in a
contextual-collapse certificate. Owners, source formulas, protection
tables, variable data and program numerals remain fixed. Logical
certificates are renamed along the same finite label map. Every
reflected premise is a generator in $\cut a\alpha$, with conclusion
$\eta$.
\end{proof}

\statement{Lemma 8 (prefixing a derived sentence)}{If
$\cut a\nu\vdash\eta$ and $K_a^\nu\eta$ is globally admissible,
then $\Lambda_a\vdash K_a^\nu\eta$ by a proof whose support is at
most $\nu$. It consequently belongs to every later cut.}

\begin{proof}
Take a finite certificate for $\eta$. Prefix its generated premises at
$\nu$ and use own validity on the certificate's valid implication.
Repeated same-label deduction gives $K_a^\nu\eta$. All body support
was strictly below $\nu$. This also prefixes derived main images after
their markers have been discharged.
\end{proof}

Two coherent lifts of one source formula on a common $\le_1$-chain,
with the same live/ghost modes, are equivalent in every cut containing
their support. To see this, compare each label with the pointwise
minimum of the two assignments. These minima preserve strict nesting:
if each parent's label exceeds its child's in both tables, the same
is true of their minima. Change children before parents using
contextual collapse, keeping every context admissible. This observation
lets a proof choose convenient canonical depth labels while preserving
the prescribed target.

\statement{Lemma 9 (guarded fixed-label transport)}{Let $\beta'$ be a
$Z$-free source-shared lift of $\beta$ with support below $\lambda$.
Assume a witnessed $G_k(\lambda)$, $\rk(\beta)<k$, and that all target
support points are $\le_1$-related to $\lambda$. Then, for every
numerical assignment $s$,
\begin{equation}\label{eq:fixed-transport}
 \cut a\lambda\vdash(\beta')^s
       \quad\Longleftrightarrow\quad S_a\vdash\beta^s.
\end{equation}
The same fixed label, mode table and target table work for every term
variant of the prescribed source skeleton and every assignment.}

\begin{proof}
Forward transport is Lemma 6. For reverse transport choose a finite
source certificate for $\beta^s$. Protect exactly the external argument
shapes of $\beta$, including their hereditary closure, and apply
$\pi_H$ to the certificate. Its target stays $Z$-free and unchanged.
Every live external argument in a transformed generating premise has
one of these protected shapes. It is a proper subexpression shape of
$\beta$, hence has rank less than $\rk(\beta)<k$. External arguments
of other shapes become ghosts, regardless of how deeply they occur.

Count active nesting through all heads, including unindexed and ghost
heads. The finite transformed certificate has some finite maximum depth.
Use the room property to choose enough whole-copy-separated
$G_{k-1}$ copies $Y_0,\ldots,Y_N$ below $\lambda$, above the fixed
target support. Put $c_t=\min Y_t$. Every copied and target point is
$\le_1$-related to $\lambda$. Connectedness therefore makes their
ordered union a single $\le_1$-chain. Label an active occurrence of
depth $t$ by $c_t$. A guard of rank $r\le k-1$ uses the first $r+1$
points of $Y_t$.

We verify each premise type. PA and Assigned Validity commute with the
logical maps. Own validity prefixes the lifted valid argument. Same-label
deduction can be used at canonical conclusion labels, and its result
can be adjusted along the common chain by contextual collapse.
Introspection uses successive depth labels. Reflection becomes live,
ghost or indexed reflection as appropriate. A source main generator
has exactly its admitted protected image, with the original source
numeral. All its live external arguments have the rank bounds just
proved, so its guards fit; discharge its holder-local markers. For an
own prefix, the proof of its body, including every point of every
earlier guard copy, precedes the next depth base. Lemma 8 supplies
the required prefix. Whole-copy separation is what justifies this last
step; separation of the bases alone would not do so.

Translate the finite valid implication. Its canonical target is
equivalent to the prescribed $\beta'$ on the common chain, so we
obtain the prescribed target below $\lambda$. The target and cutoff
have never moved. Different $s$ can require source proofs of different
depths and hence different numbers of room copies. The room property
allows every finite number; no uniform depth bound is required.

Finally, a term variant has the same protection classes, ranks and
target labels. Apply the transformation to its actual finite source
certificate, preserving its own source program numerals. This gives
the statement for all variants, not only for the initially written
formula.
\end{proof}

\statement{Lemma 10 (cofinal recovery)}{For every globally admissible,
source-shared, $Z$-free lift $\eta$ whose labels are on $C$,
\begin{equation}\label{eq:cofinal}
 \Lambda_a\vdash\eta\quad\Longleftrightarrow\quad
 S_a\vdash\eps\eta.
\end{equation}
At a head labelled $\nu\in C$, its argument's cut test already equals
the full source test below that same $\nu$, for all assignments and
term variants.}

\begin{proof}
Place the preceding finite construction on a sufficiently long part
of $C$. All positive nesting and guard requirements fit, and the
canonical target is equivalent to the given one. To identify a test
below a prescribed $\nu$, first obtain the argument proof below a
later $C$-point $\kappa$. Since $\nu\le_1\kappa$ and the argument
support is below $\nu$, Lemma 7 collapses it to $\cut a\nu$.
Only finite positive patterns are placed on $C$; arbitrary negative
diagrams are reflected in the ambient label structure instead.
\end{proof}

\section{The staged model and fully sound reconstruction}\label{sec:masks}

The staged model $M$ has standard arithmetic. Its raw tests are:
\begin{itemize}
\item at a live unindexed head $K_d$, the original predicate
$P_d((\eps\eta)^s)$ on its argument $\eta$;
\item at an active head $K_a^\nu$, the cut predicate
$\cut a\nu\vdash\eta^s$;
\item at $Z_e$, the empty predicate.
\end{itemize}
Every test is filtered by the staged truth of its argument as in
\eqref{eq:filtered}. Unmentioned original heads will have empty tests
except for those explicitly retained in full. The unindexed test uses
\emph{source erasure}; it does not compile the formula visibly
containing ghost symbols.

Lemma 4 and the commutation identities give an admissible weak-modal
structure and full modal PA before any soundness theorem is invoked.
Reflection is automatic, also at ghost heads.

\statement{Lemma 11 (truth of contextual collapse)}{Every admitted
contextual-collapse generator is true in $M$ before stage soundness
is assumed.}

\begin{proof}
At the changed atom, the two cut predicates agree by Lemma 7, and the
argument is unchanged. Through a live unindexed head, source erasure
makes the raw query identical. Through a ghost head, both tests are
false. Through an indexed outer head, the smaller-context equivalence
is itself a contextual-collapse generator for that holder below the
outer label. It identifies the two raw proof tests there. Induction
through the smaller context identifies their truth conjuncts. The whole
support of that equivalence is below the outer label by global nesting.
Boolean and quantified contexts follow by induction, at every assignment.
\end{proof}

We next reconstruct a genuine original-language filtered structure
from one main recorded argument $l\pi_H\phi$. This is the mechanism
that turns external agents into parameters we can safely use in a
lower genericity call.

Let $D$ be a set of nonactive, nonexternal names retained in full. In
the extension application it is $\{*\}$ together with a whole lower
ideal. On original closed formulas define $Q$ as follows:
\begin{enumerate}
\item Keep $Q_d=P_d$ in full for $d\in D$.
\item For active $a$, retain just its source argument shapes occurring
in $\phi$. On a sentence of such a shape, test its induced recorded
lift in the cut below the assigned label; reject all other shapes.
\item For $e\in E$, retain $P_e$ only on the argument shapes whose
$e$-head is live in the main argument. Reject ghost-selected and other
shapes.
\item Every other raw predicate is empty.
\end{enumerate}
The source-shared label and mode tables make the second and third
clauses unambiguous. The retained classes include \emph{all} term
variants and numerical instances of a shape. They do not keep merely
finitely many accepted sentences. Put $N=F_Q$.

\statement{Lemma 12 (reconstruction)}{For the retained templates and
all their term variants and assignments,
\begin{equation}\label{eq:reconstruction}
 N\models\theta[s]\quad\Longleftrightarrow\quad
 M\models l\pi_H\theta[s].
\end{equation}
Suppose in addition that the following three conditions hold:
\begin{enumerate}
\item[(S)] Every used active cut is sound in $M$.
\item[(E)] Every old $P_e$, $e\in E$, is fully sound in the original
filtered structure $F$.
\item[(X)] Under every retained live external argument, all active
tests equal their full original source tests for every term variant
and numerical assignment.
\end{enumerate}
Then every changed predicate $Q$ outside $D$ is fully sound in $N$.
The reconstructed family is admissible and recursively enumerable.}

\begin{proof}
First prove \eqref{eq:reconstruction} by formula induction, without
assuming soundness. At an active atom the raw tests agree by the
definition of $Q$. At a fully retained head they agree by erasure.
At a live external head its old raw test has been retained. At a
ghost-selected head both the original restricted test and the auxiliary
test are empty. Other named heads are empty in both structures.
Induction identifies the truth conjuncts, the Boolean connectives and
all quantified instances. Keeping all variants is essential here.

For full soundness, an active acceptance is true in $M$ by (S), and
then true in $N$ by agreement. An old external acceptance $\chi$ is
true in $F$ by (E). Heredity makes its recorded subtree $Z$-free.
Condition (X) identifies all active raw tests in that subtree; all
other occurring live tests are the same original tests. Formula
induction therefore transfers $F$-truth of $\chi$ to $M$-truth of its
lift, and \eqref{eq:reconstruction} transfers it to $N$. All remaining
predicates accept nothing. Thus no unrestricted old external predicate
is left behind with unchecked assertions in $N$.

Finite shape recognition and recorded lifting are recursive. Cut
consequence and the old predicates are recursively enumerable, so
their restrictions are recursively enumerable and alphabetically
invariant. When $D$ is a downward ideal, enumerate its full predicates
by positive finite path witnesses and dovetail their theorem
enumerations. There is no need to decide the complement of $D$.
The stipulated disjointness makes this enumeration compatible with
the finitely many changed predicates.
\end{proof}

The guard rule supplies (X). Inside a retained $K_e\chi$, an active
$K_b^\nu\beta$ has a $G_{\rk(\eps\chi)}(\nu)$ guard, with coherent
argument support below $\nu$. Its argument rank is smaller than that
guard rank. Lemma 9 therefore identifies its test with the original
source test at this fixed label. The identification holds for every
variant and assignment. Lemma 10 supplies the same identification for
$Z$-free cofinal lifts.

\subsection*{Reducing an infinite external policy to a finite one}

There is a useful preliminary restriction that does not involve
labels. For a fixed theorem $\theta$, keep the full raw predicates
that an application requires to remain full, and restrict every
outside predicate to its argument shapes occurring in $\theta$.
Hereditary closure is automatic when all subformula shapes are
included. Formula induction gives agreement with the original
structure on every retained argument variant, at all assignments.
If the old outside predicates were fully sound, their restrictions
remain fully sound. Only finitely many outside names can now be
nonempty. This allows the genericity proofs to work with finite $E$
while their conclusions apply to arbitrary admissible external families.

\section{The common theory is generic over sound externals}\label{sec:core}

The word generic here has a specific semantic meaning: the core
remains sound when other agents are replaced by any admissible raw
predicates already fully sound in the same filtered structure.
It does not assert forcing genericity.

\statement{Lemma 13 (core genericity)}{Let $F$ retain the full raw
predicate $\Th(U)$ at $*$. If every original external predicate is
admissible and fully sound in $F$, then $F\models U$.}

\begin{proof}
First assume only finitely many external predicates, with names in
$E$, are nonempty. Form the core presentation with active set $\{*\}$.
Its staged model uses the old external predicates on source erasures
and indexed core cuts; ghost and other external predicates are empty.

Verify generated sentences in increasing order of whole support.
Start with empty support. At a label $\nu$, every smaller cut is sound
because all its generated premises have already been verified.
PA and Assigned Validity are true by model admissibility. Own
validity holds because a valid argument of support below its head is
an assigned-validity consequence in the corresponding cut. Own
deduction uses closure of a single cut and the truth conjunct.
Introspection uses Lemma 8 for raw membership at its later label,
together with truth of its earlier atom. Reflection is built into the
filter, and contextual collapse holds by Lemma 11. Markers are valid
and belong to the indicated holder cuts.

Consider a main mechanicalness image. If its target is an indexed
core head $K_*^\nu\chi$, the argument support is below $\nu$.
Soundness of that smaller cut makes the truth filter redundant in
its column. The column is exactly its recursively enumerable cut
proof column. This remains true for each assignment of the other
parameters, so the existential index exists at every such assignment.

If the target is a live old $K_e\chi$, the guard rule includes this
\emph{outer} head. It supplies (X) throughout $\chi$ by Lemma 9.
Heredity makes $\chi$ $Z$-free. Old soundness in $F$ and equality of
the tests on its subtree transfer every old accepted argument to a
true staged argument. Equivalently, apply Lemma 12 to the whole
recorded modal atom, retaining its outer live shape as well, with
$D=\varnothing$. Its staged column is thus the recursively enumerable
old predicate column. A ghost or empty external target has the empty
column. This verifies all main mechanicalness images. No particular
codes occur in the core presentation.

Finally, for an own prefix at $\nu$, the body's whole support is
below $\nu$. Its truth has already been established, and it is a
generated premise in that cut. The raw test and truth conjunct
therefore both hold. Repeated prefixes are handled in the same
support order. Padding makes every invoked guard point part of the
support used in this induction. This proves all stage generators and
hence their consequences sound.

For any $U$-theorem take a $Z$-free lift on $C$. It is a stage
consequence by Lemma 10 and so true in the staged model. The same
lemma identifies every indexed raw test with its full $U$ test;
unindexed tests agree by erasure. Formula induction returns the
theorem's truth to $F$. This proves the finite-policy case without
an assumption that $U$ was already sound.

For an arbitrary external family, fix a $U$-theorem and apply the
finite shape restriction above, retaining full $\Th(U)$. Its outside
predicates remain fully sound, and its truth agrees with that in
$F$. The finite-policy case and transfer give the result.
\end{proof}

\section{Genericity for whole downward ideals}\label{sec:ideals}

For a finite nonempty $R^+$-antichain $J$, let
\[
 D_R(J)=\{i:i=a\text{ or }iR^+a\text{ for some }a\in J\}.
\]
This ideal can be infinite. It is recursively enumerable by positive
path witnesses. For $i\in\N$, define the metatheoretic rank
\begin{equation}\label{eq:rho}
 \rho(i)=\sup_{jRi}(\rho(j)+1).
\end{equation}
Well-foundedness defines $\rho$, and $iR^+a$ implies $\rho(i)<\rho(a)$.
No algorithm for these ranks will be used.

\statement{Lemma 14 (whole-ideal genericity)}{Let $F$ retain the full
raw predicates $\Th(U)$ and $\Th(T_i)$ for every $i\in D_R(J)$.
At all other original names let the predicates be admissible and
fully sound in $F$. Then $F$ models $U$ and every $T_i$ for
$i\in D_R(J)$.}

\begin{proof}
Induct on $h(J)=\max_{a\in J}\rho(a)$, for all finite nonempty
antichains. The induction assertion includes every member of the
ideal, not only the frontier $J$. The finite-policy reduction permits
us to assume that only finitely many outside predicates, with names
in $E$, are nonempty. Retain the core and the \emph{entire} ideal
unchanged in that reduction.

Form the joint extension presentations with active set $J$. The
staged model keeps the full core and full strictly lower original
predicates, uses cut predicates at the labelled $J$-heads, and the
old $E$-predicates on source erasure. Other outsiders and ghosts are
empty. Verify all holder generators simultaneously in increasing
whole support, first empty support and then successive labels.
This is the inner induction. The structural verification is exactly
the one in Lemma 13: it uses own cut closure, local prefixes,
filtered reflection, modal PA and already established collapse truth.
There are no mechanicalness main generators here. We must verify
the particular codes.

Fix a main code image with source $\phi$ and owner $w$, where $w=*$
or $wRa$ for its active holder $a$. Collect all original names in
$\phi$ that lie strictly below $J$, and also $w$ if it is original.
Let $A'$ be the $R^+$-maximal members of this finite collection, and
retain
\[
 D=\{*\}\cup D_R(A'),
 \qquad D_R(\varnothing)=\varnothing.
\]
Every member of $A'$ is strictly below a member of $J$, so if
$A'\ne\varnothing$,
\begin{equation}\label{eq:rank-decrease}
 \max_{b\in A'}\rho(b)<\max_{a\in J}\rho(a).
\end{equation}
This is a finite semantic selection for the proof. It is not a
selector in the source or stage enumeration. Every lower name
occurring in $\phi$ is covered by the retained ideal. Lower
predicates not mentioned there and outside $D$ may be emptied for
this local reconstruction.

Apply Lemma 12 to the recorded argument. Condition (S) holds because
all active tests in this main generator use cuts below its current
support maximum. Condition (E) is the old outside soundness
hypothesis. Every retained live external argument is hereditary
and $Z$-free; its guards supply (X) at fixed labels by Lemma 9.
Consequently the reconstructed original structure $N$ retains the
full core and all the lower theories in $D$, and all predicates
outside $D$ are fully sound there. The source argument is true in
$N$ exactly when its recorded lift is true in the staged model.

If $A'$ is nonempty, the outer induction hypothesis applied to
$A'$ makes $U$ and the whole retained lower ideal sound in $N$.
It includes the code owner $w$, even when $w$ is not maximal in
$A'$. If $A'$ is empty, $w$ must be the core, and Lemma 13 makes
$U$ sound in $N$ because every original predicate there is a sound
mask. Thus, for every $n$, an original $w$-proof of $\phi(\bar n)$
implies truth of its recorded lift at $n$ in the staged model.
The code head's raw query is precisely this full original proof
test, by source erasure. Equations~\eqref{eq:d} and \eqref{eq:c}
therefore verify both directions of the code: the truth filter
disappears on this argument, and the resulting test equals the
displayed program's enumeration column.

The same reasoning covers empty-support codes. There are then no
active cut or transport requirements; the lower ideal or core call
verifies their source arguments. Padding is discharged by true
markers. Own prefixes are true because their whole bodies precede
their new labels and are already true and provable at the lower
cuts. This completes the inner support induction. All active
stage consequences and cut predicates are now sound.

It remains to recover the \emph{whole} current ideal. A theorem of
$T_a$, $a\in J$, has a $Z$-free cofinal lift provable in its stage
by Lemma 10. Stage soundness makes it true. Cofinal recovery
identifies its active tests with the full original tests; every
unindexed query was already the original query on erasure.
Formula induction returns its truth to $F$. Thus all frontier
theories are sound in $F$.

Now take $i\in D_R(J)$ with $i\notin J$ and a $T_i$-theorem $\theta$.
Choose a coherent $Z$-free lift on $C$. Collect $i$ and all original
names in $\theta$ strictly below $J$, and take their finite maximal
antichain $A'$. It satisfies \eqref{eq:rank-decrease}. Reconstruct
this lift retaining $\{*\}\cup D_R(A')$ in full. Stage soundness
supplies (S), cofinal recovery supplies (X), and old outside
soundness supplies (E). The masks outside the retained ideal are
therefore fully sound. The lower induction hypothesis makes the
unchanged full $T_i$-theorem true in $N$. Reconstruction and
cofinal agreement transfer its truth to $F$. This argument works
separately for every $i$ and every theorem, even when the current
ideal is infinite.

All original predicates, inside and outside the ideal, are now
fully sound in $F$. Apply Lemma 13 at this point to obtain
$F\models U$. This closes the outer induction and proves the
finite-policy assertion. The preliminary finite shape restriction
then gives the assertion for arbitrary outside families.
\end{proof}

The proof order prevents two possible circles. Core genericity was
proved independently of extension soundness. In an extension code
verification, the outer call retains only a strictly lower original
ideal, while the inner call uses smaller active cuts. Superior and
incomparable side agents become fully sound restricted parameters;
the proof never pretends that they are subordinates.

\section{The actual family and its known schemas}\label{sec:finish}

Let $F_{\rm actual}$ be the filtered structure with full raw
$\Th(U)$ and every full raw $\Th(T_i)$. At this point we have not
assumed any of these predicates sound.

\statement{Lemma 15 (actual-family soundness)}{$F_{\rm actual}$
models $U$ and every $T_i$. Its filters are redundant.}

\begin{proof}
Fix a $T_i$-theorem $\theta$. Collect $i$ and all original modal
names occurring in $\theta$, and let $J$ be the $R^+$-maximal
members of this finite set. It is a nonempty antichain and its ideal
contains every collected name. Form $F_J$ retaining the full core
and full ideal predicates, with empty outsiders. Empty predicates
are fully sound without a hypothesis about $F_{\rm actual}$.
Lemma 14 therefore gives $F_J\models\theta$.

Every raw query occurring in $\theta$ has been retained. Formula
induction gives $F_J/F_{\rm actual}$ agreement on $\theta$, so
$F_{\rm actual}\models\theta$. Arithmetic program numerals are
unchanged, even if they encode formulas naming other agents. We
have changed semantic predicates, not theories or enumerations.
Since $i,\theta$ were arbitrary, all original theory predicates are
fully sound. Lemma 13 now gives core soundness in
$F_{\rm actual}$. Full soundness removes every truth conjunct in
\eqref{eq:filtered}.
\end{proof}

\statement{Lemma 16 (effective core inheritance)}{For every original
$i$ and closed sentence $\sigma$, a $U$-proof of $\sigma$ yields
effectively a $T_i$-proof of $\sigma$ and every finite own prefix
of $\sigma$.}

\begin{proof}
Put $\psi(x)=\sigma\wedge x=x$. A $U$-proof of $\sigma$ gives one
of $\psi(\bar0)$. By \eqref{eq:d}, the concrete $\Sigma_1$
enumeration assertion $0\in W_{d(\psi)}$ is true. Its finite
computation trace supplies a PA proof. The particular core code
in $T_i$ then gives $K_*(\sigma\wedge0=0)$; core reflection gives
$\sigma$. Use Lemma 3 to obtain its own prefixes.

The numerical instantiation of the universal code requires care
in the weak logic. Use the valid variable-instantiation implication
from $\forall x\,A(x)$ to $A(y)$ with fresh $y$, then take its
Assigned Validity instance at $y=0$. This yields the required
numerical instance without an unrestricted arbitrary-term
substitution rule inside $K_*$. The finite trace is external
evidence for this one computation, not an internal proof of the
compiler's general correctness.
\end{proof}

We can now verify the three assertions of Theorem 1. Full mutual
truthfulness is a generator of every $T_i$, and its own prefix is
also a theorem. Actual-family soundness makes that prefix true,
giving the first schema. Modal PA and the structural generators
and their prefixes similarly give self-doubting Epistemic Arithmetic.

For $jRi$, the particular sentence
\[
 \forall x(K_j\phi\leftrightarrow x\in W_{c(q,j,\phi)})
\]
and its $K_i$-prefix are $T_i$-theorems. In the actual consequence
model, assign $e=c(q,j,\phi)$ to obtain
\[
 \exists e\,K_i\forall x(K_j\phi\leftrightarrow x\in W_e).
\]
This is numeral-assignment semantics, so no illegitimate modal
term-generalization is used. Finally every $\MC_j(\phi)$ is a
closed generator of $U$. Lemma 16 gives it and its own $K_i$-prefix
in every $T_i$, hence
\[
 K_i\exists e\,\forall x(K_j\phi\leftrightarrow x\in W_e).
\]
The common core also contains parameter versions with universal
closure outside the existential, and the same transfer applies to
their closed sentences. All these formulas use the same fixed family.

\subsection*{Removing the temporary core}

For each original $i$ enumerate
\begin{equation}\label{eq:reduct}
 T_i^{\rm orig}=\{\sigma\in\mathcal L:T_i\vdash\sigma\}.
\end{equation}
This is uniformly recursively enumerable, since original-language
syntax is decidable. On original sentences, consequence from
$T_i^{\rm orig}$ equals consequence from $T_i$. One implication
composes the finite $T_i$-proofs of the premises. The other follows
because a $T_i$-theorem in the original language is itself a premise
in \eqref{eq:reduct}. Original-language validity is unchanged by
the expansion: every original weak-modal structure can be extended
by an arbitrary admissible interpretation for the extra operator.

The reduct of the sound actual structure therefore has exactly the
knowledge predicates of $T_i^{\rm orig}$. The required original EA,
truthfulness, code and mechanicalness schemas survive. The code
numerals still enumerate the right original argument columns,
because those consequence columns have just been shown identical.
The auxiliary ghosts and labels were never original axioms. Deleting
$K_*$ by this reduct leaves precisely $K_0,K_1,\ldots$. This proves
Theorem 1 once the ordinal input below is supplied. \qed

\section{Extracting the effective ordinal structure}\label{sec:ordinals}

The modal argument uses finite effective label procedures, not just
the existence of large ordinals. This section supplies one recursive
domain and total algorithms for both resemblance relations. It also
accounts for the partial conversions occurring inside the tracking
algorithms. We use the relative term procedures and structural
theorems as published mathematical inputs in Wilken
\cite{Wilken}; their numbering below refers to arXiv version 5.
They are not additional hypotheses on $R$ or on the modal family.

\subsection{Pure resemblance, reflection, and separated room}

Wilken's pure structure is
\[
 (\mathrm{Ord};<,\le_1,\le_2),
 \qquad
 \alpha\le_i\beta\ \Longleftrightarrow\
 (\alpha;<,\le_1,\le_2)\prec_{\Sigma_i}
 (\beta;<,\le_1,\le_2),
\]
with the relations defined simultaneously by induction on the upper
endpoint. Thus $\le_1$ already reflects in the language containing
$\le_2$. Ordinal addition and collapsing functions below are notation
machinery; they are not added to the reflected relational language.
We do not identify this pure structure with an arithmetic-expanded
or modified finite-covering relation.

Definition 1.5 constructs the endpoints $\up_r$. Its finite levels
are covered by the relative systems $T_{\up_r}$, starting with
$T_1$ at level zero. Theorem 1.8 gives the actual cofinal chain
$\up_2<_2\up_3<_2\cdots$ below $\up_\omega$. The upper endpoint
$\up_\omega$ is excluded. In particular, $\up_1$ has no strict
$<_2$ successor and is not the initial chain point.

For $\alpha<_1\beta$, encode the full finite diagram of a finite
support in one existential formula, with parameters for the part
already below $\alpha$. Include order and both positive and
negative resemblance facts, and the desired lower bound.
$\Sigma_1$-reflection supplies an exact copy below $\alpha$
fixing those parameters; the equality case is the identity.
This is the finite reflection used in Lemma 7. Proposition 1.6
and the following interval property supply connectedness.

For the stronger room statement, suppose
\[
 \lambda=g_0<_2\cdots<_2g_k=\zeta,\qquad k\ge1,
\]
and put $Y=\{g_0,\ldots,g_{k-1}\}$. Lemma 1.7(2) reflects $Y$
below $\lambda$, above any prescribed $u<\lambda$ and fixed finite
$X<\lambda$, with every copied point $<_1$-related to $\lambda$.
Every original point of $Y$ was $<_1$-related to $\zeta$, so the
lemma applies to the whole copy, including its last point. Repeat
with the next bound equal to the maximum of the preceding copy.
This gives $X<Y_0<\cdots<Y_{m-1}<\lambda$ with whole-copy
separation. If the fixed goal points are also $\le_1\lambda$,
connectedness makes their union with all copies one coherent
$\le_1$-chain. Initial subchains of each copy give every required
$G_r$ with $r\le k-1$. These are exactly the room quantifiers in
Lemma 9. Once the relation algorithms are available, a finite tuple
search finds such copies; their existence proves search termination.

\subsection{One public recursive domain}

Set
\[
 B_0=[0,\up_1),\qquad B_r=[\up_r,\up_{r+1})\quad(r\ge1).
\]
A public code is a normalized finite $T_1$ term in $B_0$, or a
finite $T_{\up_r}$ term in $B_r$ whose occurring parameters below
$\up_r$ are fully written public trees from earlier blocks. The
endpoint $\up_r$ has a finite tag and its native base term
$\vartheta^{\up_r}(0)$; when used in a later block it is an earlier
parameter. This introduces no self-reference.

Definition 2.17, Theorem 2.18, and the paragraph after Corollary 2.19
give one uniform relative grammar and comparison procedure, with
supplied parameter data. The largest written collapsing index
selects a sufficient finite fragment (Definition 2.21). Finite
subterm sets and the comparison criterion of Lemma 2.26 determine
the normal-form and countable-sort tests. There is no search
through semantic hulls. The formal $\Omega_j$ symbols are finite
grammar sorts, not input oracles for uncountable cardinals.

Calls to parameter data strictly decrease the finite block tag.
Within a block, the relative procedures recurse on finite syntax.
Carry additive and multiplicative normal forms, ordinal arithmetic,
and principal-number classification with the earlier-block data;
a bare recursive ordering of parameters would not suffice. Epsilon
membership uses Lemma 2.29, fixed-point comparisons use Lemma 2.30,
and logarithm and exponentiation use Lemmas 2.41 and 2.51. The
epsilon endpoints are closed under the finite arithmetic operations
used here. Results below a base become earlier parameter codes;
other results have the normalized code in their containing block.
An exact division in a guaranteed branch can alternatively be found
by enumerating public codes and testing the witnessed multiplication
equation. Existence at its caller makes that search terminate.

If hereditary spellings have duplicates, first allow recursive raw
codes with recursive ordinal equality. Keep exactly those raw codes
having no numerically smaller equal valid raw code. This is a finite
recursive check. Each equality class has a least numerical code;
converting a computed output checks only codes up to that output.
Thus uniqueness need not be a separate assumption.

Coverage of each block follows from Theorem 2.18 and Definition 1.5.
The same procedure works for every finite block tag, so this is one
recursive presentation of $[0,\up_\omega)$, not a nonuniform union
of separately recursive orders. Arithmetic and comparison are
constructed before tracking or resemblance is used. On this domain
the limit-segmentation parameter in the source is zero, and no
nonzero limit-index endpoint has to be coded.

\subsection{Typed local bases and the conversion invariant}

Tracking procedures also use temporary proper epsilon bases.
Retain the public code of each ordinal while carrying, when needed,
a native $T_b$ representation with publicly coded parameters below
$b$. A native term is not silently a fixed-endpoint public term.
Every conversion has an explicitly checked source context and bound.

Let $v$ be the public block base, and let $a$ be a proper epsilon
in that block. Use its relative successor bound $a^+$ from
Convention 2.28. For countable inputs, the star-subterm estimate
of Lemma 2.26 and Remark 2.27 makes
\begin{equation}\label{eq:conversion}
 \xi<a^+\quad\Longrightarrow\quad
 \text{the public/native conversions between the $v$ and $a$
 contexts are defined on $\xi$}.
\end{equation}
These are the partial translations on page 17. Temporary high-sort
expressions instead carry the actual star-subterm restriction in
Definition 2.45. The routine $\iota_{v,a}$ acts on admitted parts
of the collapsing argument whose star bounds are below $a$.

The important outputs obey
\begin{equation}\label{eq:lambda-mu}
 \lambda_a<a^+,
 \qquad \mu_a<a^+,
 \qquad x\le\mu_a\ \Longrightarrow\ \rho_x^a\le\lambda_a<a^+.
\end{equation}
The references are Lemmas 2.48, 3.7 and 3.14(b). Thus these native
results return to their public context. For consecutive proper
bases of a $\mu$-covered reference sequence,
$a_{i+1}\le\mu_{a_i}<a_i^+$. Lemmas 2.48(d), 3.8, 3.26 and the
nested localization bounds in Lemma 2.33 identify the suppressed
relative contexts on their overlaps.

To resolve a parameter $p$ with $a_i\le p<a_{i+1}$, use
\[
 p<a_{i+1}\le\mu_{a_i}<a_i^+
\]
to convert its existing public payload into $T_{a_i}$. The case
$p=a_i$ is the base term. Parameters below $a_i$ move to an earlier
reference position. After endpoint positions are reached, calls
decrease the public finite block tag. This is an effective decreasing
resolution order, and produces the actual native resolved syntax
of Definitions 3.39--3.41. Its natural-valued length is computable;
it is not the length of some different public normalization.

Endpoint bases have explicit clauses in Definition 3.15: their
$\lambda/\mu$ output is the next finite endpoint. Interior inputs
already have the public native block representation; equality with
the next endpoint is dispatched before a proper-base conversion.
Zero has its own $\kappa,\mathrm{dp},\mathrm{tc}$ cases. No caller
asks for $\lambda_0$ or $\mu_0$.

\subsection{Tracking sequences and finite evaluation}

Here is the procedure order and the conversion information needed by
its callers. The underlying finite operations are Wilken's published
relative procedures; the preceding invariant puts them on the one
public datatype.

Localization and fine localization process finite principal-subterm
sets (Definitions 2.32, 2.34, 2.35 and 2.38). Fixed-point tests use
Lemma 2.30 rather than semantic supremum searches. The terminating
bar operations are justified in Lemmas 2.35--2.40. The $\chi$ and
$\zeta$ data are finite normal-form operations (Definitions 3.3
and 2.42); $\lambda$ and $\mu$ use the translations with the
bounds in \eqref{eq:lambda-mu}. A $\rho$ caller has $x\le\mu_b$,
so its output has the same justified return conversion.

Tracking-sequence assignment $\mathrm{ts}$ uses finite localization
recursion for multiplicative principal inputs (Definition 3.19,
Lemma 3.21), and finite multiplicative decomposition for other
principal inputs (Definition 3.20). A temporary call
$\mathrm{ts}^{a_i}(\beta)$ has the source bound $\beta\le a_n$;
nested localization puts it below $a_i^+$ and permits its conversion.
Endpoint extension is the finite prefix procedure of Definition 3.24.
The covering procedures of Definitions 3.25--3.28 select finite
localization subsequences.

Two evaluation details matter for effectiveness. The bound
$\widehat\gamma$ in Definition 3.30 is introduced as a semantic
minimum, but is not computed by a least-ordinal search. At an actual
$h$-call in $o$, the input has $\beta\le\gamma$ or
$\beta\le\mu_\gamma$. Since $\widehat\gamma$ exceeds both, the
advertised $h$-domain follows, and the computable conversion bound
is $\beta<\gamma^+$. The evaluator can branch on the ordinary
comparisons with $\gamma$ and $\mu_\gamma$.

Also, the overlined operation in Definition 3.33 truncates factors
strictly below $\beta$ from the multiplicative normal form of
$\mu_\delta$. Its output is at most $\mu_\delta$. It is not the
ordinary product $\mu_\delta\cdot\beta$. The actual truncation
preserves the covered next entry and the decreasing native-height
bound of Lemma 3.7. Thus $\mathrm{sk}$ and $h$ are finite typed
operations. Evaluation $o$ then uses the finite sequence-length
descent of Definition 3.35 and Theorem 3.36, after the endpoint
fixed-point cases. Final sums and products are formed on public
codes. A large evaluated product is not passed through a local
partial translation merely because its factors had native types.
Lemma 3.49 and Corollary 3.50 give finite containing blocks for
these evaluations.

\subsection{Scheduling the component procedures without a cycle}

Definitions 3.37, 3.43 and 3.44 appear mutually dependent if read
as a single informal program. They have the following finite
construction order:
\begin{enumerate}
\item Define principal $\kappa$ and principal $\nu$ at indices
greater than $1$ by the already constructed evaluator $o$.
For nonempty contexts set $\nu_0=o(\vec{\alpha})$ and
$\nu_1=o(\vec{\alpha})\cdot2$ directly. At a proper last
base $b$, the respective native domains are bounded by
$\lambda_b$ and $\mu_b$, so \eqref{eq:lambda-mu} licenses conversion.
\item Define general $\kappa$ and $\mathrm{dp}$ by Definition 3.43,
with principal $\nu$ already available.
\item Define nonprincipal $\nu$ by Definition 3.44, using the
already total $\kappa$ and $\mathrm{dp}$ procedures.
\end{enumerate}
In the epsilon clause for $\mathrm{dp}$, the call is
$\nu^{\vec{\alpha}\frown b}_{\mu_b}$. The index $\mu_b$ is
additive principal, being the stated $\omega$-power or a finite
endpoint, and is greater than $1$. Thus this is a call to the
\emph{principal} evaluator from step 1, not to the later general
$\nu$ recurrence. It satisfies the last-index restriction in
Definition 3.18. No recursive cycle remains in this scheduling.

The measure for step 2 is lexicographic:
\begin{equation}\label{eq:length-measure}
 (\text{native resolved natural term length},
   \text{reference-list length}).
\end{equation}
Additive clauses use proper nonzero components. The non-epsilon
$\mathrm{dp}$ clause uses the shorter resolved logarithm or
divided-logarithm arguments of Lemma 3.42(2). The epsilon clause
appends $b$ and calls $\kappa/\mathrm{dp}$ at $\lambda_b$;
Lemma 3.42 supplies the strict native resolved-length decrease
after that append. The same-value clause removes a last context
base above the input: its fully resolved payload is unchanged and
the reference-list length decreases. Zero, one, base and endpoint
cases terminate directly.

The appended list must itself be covered. Its caller has
$b\le\lambda_a$ for the preceding base $a$, and Lemma 3.14(c)
gives $b\le\mu_a$. If $\mu_a<a$, then $\lambda_a<a^2$ and
there is no new proper epsilon greater than $a$ in that interval;
otherwise the greatest epsilon below $\lambda_a$ equals that below
$\mu_a$. Endpoint cases have the explicit next-endpoint bound.
This verifies the reference-sequence hypothesis of the native
length estimate.

Step 3 recurses on a successor predecessor or proper additive
components. Its additional $\kappa/\mathrm{dp}$ calls at
$\rho$-indices invoke procedures already made total; they do not
require a claim that every $\rho$ term is shorter than the original
$\nu$ input. Lemma 3.14(b) supplies their local $\lambda$-domain.
Equations~\eqref{eq:conversion}--\eqref{eq:lambda-mu} and the
parameter-resolution procedure provide the native representations
at each recursive call and public codes on return.

\subsection{The relative-tracking caller bound}

Definition 3.85 constructs $\mathrm{ts}[\vec{\alpha}](z)$ by
finding the largest position $k$ with
$o(\vec{\alpha}\upharpoonright k)\le z$ in a finite list.
The $k=0$ case uses the public endpoint segment. For $k>0$ it
forms the public arithmetic value
\begin{equation}\label{eq:relative-z}
 z'=a_k\cdot\bigl((1/o(\vec{\alpha}\upharpoonright k))\cdot z\bigr).
\end{equation}
The parenthesized expression denotes the source's exact ordinal
quotient, not an ordinary real reciprocal. Its existence is
Lemma 3.53. Before conversion to the native $T_{a_k}$ context we
need the actual bound $z'<a_k^+$; $z<\up_\omega$ alone would not
give it.

If $k$ is not the last position, maximality gives
$z<o(\vec{\alpha}\upharpoonright(k+1))$, and Lemma 3.49
bounds the latter by
$o(\vec{\alpha}\upharpoonright k)\cdot\widehat a_k$.
At the last position the declared domain gives
$z\le\kappa^{\vec{\alpha}}_{\lambda_{a_k}}+
\mathrm{dp}^{\vec{\alpha}}(\lambda_{a_k})$;
Corollary 3.50 bounds that expression strictly below
$o(\vec{\alpha})\cdot\widehat a_k$. In both branches,
\[
 z<o(\vec{\alpha}\upharpoonright k)\cdot\widehat a_k.
\]
The exact quotient in \eqref{eq:relative-z} is therefore below
$\widehat a_k$. This bound is multiplicative principal and strictly
greater than $a_k$. Closure under the product of two smaller
factors gives
\begin{equation}\label{eq:relative-bound}
 z'<\widehat a_k\le a_k^+.
\end{equation}
The last comparison is stated following Definition 3.30.
No computation of the semantic $\widehat a_k$ is needed: it is
used to prove a bound on the public arithmetic output.
Equation~\eqref{eq:relative-bound} licenses its actual partial
conversion. Lemmas 3.52--3.53 place the output tracking sequence
in the required relative initial segment. If $a_k$ is an endpoint,
the native public endpoint block is used directly.

\subsection{Tracking chains and total relation decisions}

Index-chain and reference data are finite scans of publicly coded
lists (Definitions 3.54--3.62 and 3.74). Reference pointers go to
earlier index pairs, and Lemmas 3.63 and 3.75 give the component
domains. Maximal extension uses the typed $\mu,\lambda$, logarithm
and $\rho$ procedures (Definitions 3.64 and 3.71). Lemma 3.67 proves
termination after contracting adjacent $\mu/\lambda$ pairs; its
endpoint clause stops rather than iterating the next-endpoint
operation indefinitely. Recognition of tracking chains uses shorter
proper initial chains and finite extension lists (Definitions
3.68--3.69). Chain evaluation is a finite fold of the already total
component procedures (Definition 3.76).

The tracking-chain assignment $\mathrm{tc}$ of Definition 3.87
uses $\mathrm{ts}$ for principal inputs and finite additive
decomposition otherwise. Every relative tracking preparation has
the bound just proved. Lemma 3.89 verifies both the chain condition
and evaluation back to the input. The greatest branching index
of Definition 4.5 recurses on proper prefixes. Consequently the
reach calculation $\operatorname{lh}$ in Corollary 4.6 is effective
on the one public domain. Its result is a public finite term or
the separate sentinel $\infty$, never an uncoded ordinal passed
to an arithmetic procedure.

The closed-interval property now gives the total decision
\begin{equation}\label{eq:le1}
 \alpha\le_1\beta\quad\Longleftrightarrow\quad
 \alpha\le\beta\ \text{and}\
 \bigl(\operatorname{lh}(\alpha)=\infty
       \ \text{or}\ \beta\le\operatorname{lh}(\alpha)\bigr).
\end{equation}
Endpoint equality is recursive on public codes;
$\operatorname{lh}(0)=0$ and
$\operatorname{lh}(\up_r)=\infty$ for $r\ge1$ are explicit cases.

Theorem 4.2(b) computes the greatest strict $<_2$ predecessor
$p$ of a nonminimal $\beta$, apart from its stated nonzero
limit-index endpoint exception. Such endpoints do not occur
below $\up_\omega$. Minimality is a finite tracking test.
If $\alpha<_2\beta$, greatestness gives $\alpha\le p$ in the
ordinary order. Since $\alpha,p\le_1\beta$, connectedness gives
$\alpha\le_1p\le_1\beta$. The pure respect property
\[
 u\le_1v\le_1w,\quad u\le_2w\quad\Longrightarrow\quad u\le_2v
\]
therefore gives $\alpha\le_2p$. Transitivity supplies the converse,
so
\begin{equation}\label{eq:le2-descent}
 \alpha<_2\beta\quad\Longleftrightarrow\quad\alpha\le_2p.
\end{equation}
Test equality first, then ordinary order and minimality, and otherwise
replace $\beta$ by its computed greatest predecessor. Each iteration
is a total finite procedure and strictly decreases the represented
ordinal. Well-foundedness proves termination; a recursive running-time
bound is unnecessary. Zero cannot conceal a real predecessor, because
it has no strict $\le_2$ successor. As endpoint checks,
$\up_1$ is minimal without strict successors, $\up_2$ is minimal
with strict successors, and $\up_n$ has predecessor $\up_{n-1}$
for $n\ge3$.

\statement{Lemma 17 (effective label package)}{There is one recursive
presentation of $[0,\up_\omega)$ with recursive $<,\le_1,\le_2$
and all four properties in Section~\ref{sec:room}.}

\begin{proof}
The public grammar, normalization and arithmetic precede tracking.
Every temporary conversion has its public payload and the caller
bounds above. Finite parameter resolution, the evaluation descent,
the scheduled component recursion and finite chain operations return
public codes and terminate. Equations~\eqref{eq:le1} and
\eqref{eq:le2-descent} give total relation algorithms. Abstract
reflection, separated room and the cofinal chain were established
in the first subsection. Finite searches for guard copies use these
relation algorithms, with termination supplied by room. No operation
returns $\up_\omega$ as a member code; the only unbounded reach
is the explicitly separate $\infty$ sentinel. This supplies exactly
the package used in the modal proof.
\end{proof}

\section*{Scope of the argument}

The proof above covers the complete 3-1-3 statement: arbitrary
recursively enumerable well-founded $R$, unrestricted modal formulas,
one uniform family, all three known schemas, and the original-language
reduct. Its nontrivial published dependencies are Alexander's weak
modal completeness and arithmetic semantics, the ordinary recursion
and parameter theorems, and Wilken's relative ordinal procedures and
pure resemblance theorems. The uniform extraction of those procedures
is included rather than left as an extra effective-ordinal assumption.

This manuscript expands a complete proposed hand proof whose recorded
modal and ordinal checks have been completed. Those checks do not
constitute a proof-assistant certificate or external peer review.
The separate Lean development still has unfinished dependencies and
proofs. In particular, no unconditional Lean theorem for full 3-1-3
is asserted here. The source correspondence and exact document-build
checks are recorded alongside this paper.

\begin{thebibliography}{9}
\bibitem{Alexander}
Samuel Allen Alexander.
\emph{The Theory of Several Knowing Machines}.
Ph.D. dissertation, The Ohio State University, 2013.
\href{https://etd.ohiolink.edu/acprod/odb_etd/ws/send_file/send?accession=osu1366362999\&disposition=inline}{OhioLINK dissertation record and text}.

\bibitem{Wilken}
Gunnar Wilken.
Pure $\Sigma_2$-elementarity beyond the core.
\emph{Annals of Pure and Applied Logic} 172 (2021), article 103001.
\href{https://doi.org/10.1016/j.apal.2021.103001}{doi:10.1016/j.apal.2021.103001}.
Source numbering follows
\href{https://arxiv.org/abs/1710.01870v5}{arXiv:1710.01870v5}.
\end{thebibliography}

\end{document}
